\documentclass[a4paper]{article}

\usepackage[utf8]{inputenc}
\usepackage{amsmath}
\usepackage{graphicx}
\usepackage{xcolor}

\setlength{\parindent}{0pt}
\setlength{\parskip}{0.5em}

\definecolor{thmcolor}{RGB}{220,235,255}
\definecolor{defcolor}{RGB}{232,247,228}
\definecolor{excolor}{RGB}{255,244,220}

\providecommand{\mathbb}[1]{\mathbf{#1}}
\providecommand{\mathcal}[1]{\mathit{#1}}
\newcommand{\cref}[1]{\ref{#1}}
\newcommand{\toprule}{\hline}
\newcommand{\midrule}{\hline}
\newcommand{\bottomrule}{\hline}
\newcommand{\href}[2]{#2}
\newcommand{\url}[1]{\texttt{#1}}
\renewcommand{\labelitemi}{-}
\renewcommand{\labelitemii}{--}

\newcommand{\boxheading}[1]{%
  \par\smallskip
  \noindent\textbf{\ifx&#1&Note\else#1\fi}\par
  \smallskip
}

% Theorem environment
\newtheorem{theorem}{Теорема}[section]
\newenvironment{thmbox}[1][]{%
  \begin{quote}\color{blue!60!black}\boxheading{#1}\normalcolor
}{\end{quote}}

% Definition environment
\newtheorem{definition}{Определение}[section]
\newenvironment{defbox}[1][]{%
  \begin{quote}\color{green!40!black}\boxheading{#1}\normalcolor
}{\end{quote}}

% Lemma environment
\newtheorem{lemma}{Лемма}[section]
\newenvironment{lembox}[1][]{%
  \begin{quote}\color{orange!70!black}\boxheading{#1}\normalcolor
}{\end{quote}}

% Proposition environment
\newtheorem{proposition}{Утверждение}[section]
\newenvironment{propbox}[1][]{%
  \begin{quote}\color{violet!70!black}\boxheading{#1}\normalcolor
}{\end{quote}}

% Corollary environment
\newtheorem{corollary}{Следствие}[section]
\newenvironment{corbox}[1][]{%
  \begin{quote}\color{cyan!50!black}\boxheading{#1}\normalcolor
}{\end{quote}}

% Example environment
\newtheorem{example}{Пример}[section]
\newenvironment{exbox}[1][]{%
  \begin{quote}\color{brown!70!black}\boxheading{#1}\normalcolor
}{\end{quote}}

% Remark environment
\newtheorem{remark}{Замечание}[section]
\newenvironment{rembox}[1][]{%
  \begin{quote}\color{black!70}\boxheading{#1}\normalcolor
}{\end{quote}}

% Customize enumerate and itemize
% Custom table environment with caption, placement, and column spec
\newenvironment{customtable}[3][htbp]{%
  % #1 = optional: placement (default: htbp)
  % #2 = mandatory: column specification (e.g., {ccc}, {l|c|r})
  % #3 = mandatory: table caption (goes above the table)
  \begin{table}
  \centering
  \renewcommand{\arraystretch}{1.2}
  \textbf{#3}\par\smallskip
  \begin{tabular}{#2}
}{%
  \end{tabular}
  \end{table}
}

% ============================================================
% DOCUMENT STRUCTURE GUIDELINES
% ============================================================
%
% === 1. DOCUMENT STRUCTURE & FLOW ===
% - Use \section{}, \subsection{}, \subsubsection{} to mirror lecture flow.
% - Limit depth: section → subsection → subsubsection (avoid deeper).
% - Start each section with a 1--2 sentence summary of what follows.
% - Example:
%   \section{Introduction}
%   This section introduces the core concepts of vector spaces...

% === 2. CONTENT ACCURACY ===
% - Faithfully represent the lecture: definitions, theorems, examples, logic.
% - Do NOT add external content unless noted (e.g., in footnotes).
% - Paraphrase for clarity; use direct quotes sparingly.

% === 3. MATHEMATICAL TYPESETTING ===
% - Always use math mode: $x^2$, \[ \int f(x) dx \], or \begin{equation}.
% - Number equations only if referenced later (use \label{} and \cref{}).
% - Use semantic commands: \mathbb{R}, \mathcal{L}, \vec{v}, \nabla, etc.
% - Example:
%   \begin{equation}\label{eq:euler}
%   e^{i\pi} + 1 = 0
%   \end{equation}
%   As shown in \cref{eq:euler}, ...

% === 4. CUSTOM ENVIRONMENTS (USE AS INTENDED) ===
% - \begin{defbox}[Title]     → For definitions (green-tinted).
% - \begin{thmbox}[Title]     → For theorems/lemmas/propositions (blue-tinted).
% - \begin{exbox}[Title]      → For examples (beige-tinted).
% - \begin{rembox}[Title]     → For remarks/notes (gray-tinted).
%
% Rules:
% - Always include a title: \begin{thmbox}[Pythagorean Theorem]}
% - Keep content concise; avoid long paragraphs inside boxes.
% - Do NOT nest boxes.
% - For formal numbering, wrap amsthm environments inside:
%   \begin{thmbox}[Mean Value Theorem]
%   \begin{theorem}
%   Let $f$ be continuous on $[a,b]$...
%   \end{theorem}
%   \end{thmbox}

% === 5. LISTS AND ENUMERATIONS ===
% - Use itemize for unordered key points:
%   \begin{itemize}
%     \item Point A
%     \item Point B
%   \end{itemize}
% - Use enumerate for step-by-step procedures:
%   \begin{enumerate}
%     \item Step 1: Define domain.
%     \item Step 2: Compute derivative.
%   \end{enumerate}
% - Avoid nesting >2 levels.

% === 6. TABLES (USE customtable) ===
% - Syntax:
%   \begin{customtable}[placement]{column-spec}{Caption}
%   \toprule
%   Col1 & Col2 & Col3 \\
%   \midrule
%   Data & Data & Data \\
%   \bottomrule
%   \end{customtable}
% - Example:
%   \begin{customtable}[h]{lcc}{Experimental Results}
%   \toprule
%   Trial & Input & Output \\
%   \midrule
%   1 & 10 & 15 \\
%   2 & 20 & 30 \\
%   \bottomrule
%   \end{customtable}
% - Use booktabs rules (\toprule, \midrule, \bottomrule).
% - No vertical lines; align numbers by decimal if needed.

\begin{document}

\begin{center}
{\LARGE \textbf{Исследование случаев с τ(k₂) < 0 и τ(k₂) > 0 в математическом анализе}}\\[1em]
\end{center}

\textbf{Abstract.} В лекции рассматриваются два различных случая функции τ(k₂): когда τ(k₂) меньше нуля и когда τ(k₂) больше нуля. Основной целью является понимание поведения функции в зависимости от знака τ(k₂). Обсуждаются ключевые концепции, связанные с анализом функций, а также возможные применения этих случаев в математических моделях. Приводятся примеры, иллюстрирующие различия в поведении функции в каждом из случаев, что позволяет сделать выводы о влиянии знака на свойства функции и её графическое представление. Заключение подчеркивает важность различения этих случаев для дальнейшего анализа и применения в различных областях математики.

\begin{center}
{\LARGE \textbf{Исследование случаев с τ(k₂) < 0 и τ(k₂) > 0 в математическом анализе}}\\[1em]
\end{center}

\textbf{Abstract.} В лекции рассматриваются два различных случая функции τ(k₂): когда τ(k₂) меньше нуля и когда τ(k₂) больше нуля. Основной целью является понимание поведения функции в зависимости от знака τ(k₂). Обсуждаются ключевые концепции, связанные с анализом функций, а также возможные применения этих случаев в математических моделях. Приводятся примеры, иллюстрирующие различия в поведении функции в каждом из случаев, что позволяет сделать выводы о влиянии знака на свойства функции и её графическое представление. Заключение подчеркивает важность различения этих случаев для дальнейшего анализа и применения в различных областях математики.

\section{Введение}
В данном разделе мы рассматриваем функцию τ(k₂) и её поведение в зависимости от знака. Функция τ(k₂) может принимать как отрицательные, так и положительные значения, что существенно влияет на её свойства и графическое представление.

\section{Случай τ(k₂) < 0}
\subsection{Описание поведения функции}
Когда τ(k₂) меньше нуля, функция может демонстрировать определённые характеристики, такие как убывание или наличие экстремумов. Важно проанализировать, как это влияет на график функции и её поведение в окрестности критических точек.

\subsection{Примеры}
\begin{exbox}[Пример 1]
Рассмотрим функцию τ(k₂) = -k₂^2. В этом случае функция принимает отрицательные значения для всех k₂, что приводит к параболическому графику, открывающемуся вниз.
\end{exbox}

\section{Случай τ(k₂) > 0}
\subsection{Описание поведения функции}
При τ(k₂) больше нуля функция может вести себя совершенно иначе, например, демонстрируя возрастание или наличие локальных минимумов. Это также важно для понимания её графического представления.

\subsection{Примеры}
\begin{exbox}[Пример 2]
Рассмотрим функцию τ(k₂) = k₂^2. В этом случае функция принимает положительные значения для всех k₂, что приводит к параболическому графику, открывающемуся вверх.
\end{exbox}

\section{Заключение}
В заключение, различие между случаями τ(k₂) < 0 и τ(k₂) > 0 имеет важное значение для анализа функций. Понимание этих случаев позволяет более глубоко исследовать свойства функций и их применение в различных областях математики.\end{document}